\documentclass[10pt,a4paper]{amsart}
\usepackage[margin=19mm]{geometry}
\usepackage{amsmath,amssymb,mathtools}
\usepackage[scr=rsfs,cal=euler]{mathalfa}
\usepackage{xfrac}
\usepackage[unicode,hidelinks,bookmarks=false]{hyperref}
\newcommand{\ie}{i.e.\ }
\newcommand{\cf}{cf.\ }
\newcommand{\resp}{respectively\ }
\newcommand{\Subsection}{\subsection}
\providecommand{\nobreakdash}{\nobreak\hbox{-}}
\setlength{\emergencystretch}{2em}
\multlinegap0pt
\usepackage{dsfont}
\usepackage{amssymb}
\usepackage{enumerate}
\newtheorem{thm}{Theorem}
\newcommand{\sX}{{\mathcal X}}
\title{Fixed point locus of Moduli spaces of sheaves on Toric Dm stacks}
\author{Promit Kundu}
\date{Source: arXiv:2605.01730v1 (CC BY 4.0)}
\begin{document}
\maketitle

\noindent\emph{Source and licence.} Fixed point locus of Moduli spaces of sheaves on Toric Dm stacks. Version arXiv:2605.01730v1 (CC BY 4.0), distributed by arXiv under the Creative Commons Attribution 4.0 International licence, http://creativecommons.org/licenses/by/4.0/, as recorded in the arXiv licence metadata for this exact version and shown on the abstract page https://arxiv.org/abs/2605.01730v1. Full licence notice: \texttt{licenses/arxiv-2605.01730-CC-BY-4.0.txt}.\par\smallskip

\noindent\textit{Editorial scope.} Counted unit: the article's thm Formula (source0.tex lines 988--1004). The article's complete original proof of it is delivered with it (source0.tex lines 1006--1104, one \texttt{proof} environment of 99 lines). No mathematics is written, repaired or re-derived: the statement and the proof are the article's own lines, in the article's own notation, and no retained line is substituted or re-flowed.  Retained uncounted as the source-local context the proof needs: lines 924--926.  Nothing was omitted from any retained span, so the package carries no omission record.  No retained line contains a citation, so the wrapper carries no bibliography and no bibliography entry.  No retained line contains a figure environment, an image-inclusion command or drawing code, so no image is delivered and the package declares no asset dependency.  Editorial formatting only, in addition: an AMS \texttt{amsart} preamble that restates the article's own declarations and macros used by the retained lines, namely the environments \texttt{thm} on separate counters, together with the standard packages those lines need.  The retained environments print in this document as Theorem 1 in order of appearance, whereas the article prints the same items with the numbering of its own full document.  The complete native source of the article as distributed in the arXiv e-print for this exact version is bundled unchanged as \texttt{sources/199655/source0.tex}.
\begin{equation}\label{cokernel}
0\to DH\to DG_{\sigma_1}\bigoplus DG_{\sigma_2}\to X(K)\to 0.    
\end{equation}

\begin{thm}\label{Gluing Formula}
Let $\sX$ be a smooth toric DM stack given by $(N,\Sigma,\beta).$ 
%with top cones given by $\sigma_1,..,\sigma_l$ and let $(\rho^{(i)}_{1},..,\rho^{(i)}_{d})$ be the rays corresponding to top cones $\sigma_i.$ Let $\{\rho^{(i)}_{i_{1}},..,\rho^{(i)}_{i_{d-p}}\} \subset \{{\rho^{(i)}_{1},..,\rho^{(i)}_{d}}\}$ respectively $\{\rho^{(j)}_{j_{1}},..,\rho^{(j)}_{j_{d-p}}\} \subset \{{\rho^{(j)}_{1},..,\rho^{(j)}_{d}}\}$ be the rays of $\sigma_{i} \cap \sigma_{j}$ labeled in such a way that
%$\rho^{(i)}_{i_{k}}=\rho^{(j)}_{j_{k}}$ for $k=1,..,d-p.$
The category of torsion free toric sheaves is equivalent to the category of $l$-tuples $\{\hat{F}_i\}_{i=1,.,l}$ of finite stacky $S$- families on $\mathscr{U}_{i=1,.,l}$ satisfying the following equality of $K_{i,j}$ representations given at each point $Q+\sum_{i=1}^{d-p}l_i\chi_i$ where $Q\in Box(\chi_1,..,\chi_d),\ i\neq j, \ (l_1,..,l_{d-p})\in \mathbb{Z}^{d-p}$ by,

$$\bigoplus_{\begin{array}{cc}
     &  l\in DG_{\sigma_i}\\
     & Q'\in S_{\sigma_{i,ij,Q}}
\end{array}} ({}_{Q'}F_{i}(l_1,.,l_{d-p},\infty,.,\infty)_{l} \otimes\sum_{i=d-p+1}^{d}(l_{i}\eta_{i},0))\ \otimes\sum_{i=d-p+1}^{d}a_{i-d+p}\mu_{i}$$
$$=\bigoplus_{{\begin{array}{cc}
     &  l\in DG_{\sigma_j}\\
     & Q'\in S_{\sigma_{j,ij,Q}}
\end{array}}}({}_{Q'}F_{j}(l_1,.,l_{d-p},\infty,.,\infty)_{l}\otimes\sum_{i=d-p+1}^{d}\ (0,l'_i\eta'_i))\ \otimes \sum_{i=d-p+1}^{d}a'_{i-d+p}\mu_{i}\ $$
%$$+ \sum_{i=1}^{d}\overline{(0,(l'_{i}-l_{i})\eta'_{i}}).$$ %\ \sum_{i=1}^{d-p}\big ((l'_i-l_i)\mu_{i}-\sum_{r=1}^{p}t_{i,r}\mu_{d-p+r}\big ).
Similar gluing conditions hold for the morphisms.
\end{thm}

\begin{proof}
We prove the theorem for i=1 and j=2 and can be easily generalized to other top cones. In the rest of the proof $K_{i,j}$ is basically denoted as $K_{1,2}=K_{2,1}=K.$
We abbreviate $S_{\sigma_{j,ij,Q}}$ with $S_{\sigma_1,Q}$ and similarly for $\sigma_2.$ We use the same set of torus weights $(\chi_1,..,\chi_d)$ for $\mathscr{U}_{12}.$\par 
Consider $\mathscr{F}$ restricted to $Z_1$ denoted by $\mathscr{F}_1.$ Let us denote by ${}_{(q_1,..,q_d)}F_1(l_1,,l_d)$ the weight sum decomposition of the sheaf indexed by $(q_1,,q_d)\in B_{\sigma_{1}}(\mathbb{T})$ given by the basis of $\mathbb{T}-$ weights on $Z_1.$ On restricting $\mathscr{F}_1$ to $\mathbb{C}^{d-p}\times \mathbb{C^*}^p$ we obtain that for fixed $(l_1,.,l_{d-p})$ the vector spaces ${}_{(q_1,..q_d)}F_1(l_1,,.,l_d)$ stabilize for $l_{d-p+1}>>0,..,l_{d}>>0.$
Upto isomorphism of the stacky finite S-family $\hat{F_1}$ we can take the isomorphisms between the limiting spaces as identity and then denote the limit by: $${}_{(q_1,..,q_d)}F_1(l_1,..l_{d-p},\infty,..,\infty).$$ Consider the $S-$ family corresponding to the sheaf $\mathscr{F}_1$ restricted to $\mathbb{C}^{d-p}\times \mathbb{C^*}^p$ say $\hat{G_1}$ and from the description of graded tensor product one obtains:
$${}_{(q_1,.,q_d)}G_{1}(l_1,.,l_d)={}_{(q_1,.q_d)}F_{1}(l_1,.,l_{d-p},\infty,.,\infty).$$ For sufficiently large $l_{d-p+1},.,l_d$ we have $${}_{(q_1,.,q_d)}F_1(l_1,.,l_d)={}_{(q_1,.,q_d)}F_1(l_1,.,l_{d-p},\infty,.,\infty)$$ and hence the fine grading is obtained as $l\in DG_{\sigma_1}$:
$${}_{(q_1,.,q_d)}G_{1}(l_1,.,l_d)_l={}_{(q_1,.,q_d)}F_{1}(l_1,.,l_d)_l.$$ The fine grading for a fixed $(l_{d-p+1},.,l_{d})$ determines the fine grading for any collection of $(l_{d-p+1},.,l_{d})$ given by:
$${}_{(q_1,.,q_d)}G_1(l_1,.,l_d)_{l}={}_{(q_1,.,q_d)}G_1(l_1,.,l_{d-p},0,.,0)_{l-\sum_{i=d-p+1}^{d}\eta_il_i}\otimes \eta_{d-p+1}^{l_{d-p+1}}..\otimes \eta_{d}^{l_{d}}.$$
The above notation means that the character $$\eta_{d-p+1}^{l_{d-p+1}}\otimes.. \eta_{d}^{l_{d}}$$ is a function from $G_{\sigma_{1}}$ to $\mathbb{C}^{*}$ and the representation refers to the action of an element of $G_{\sigma_{1}}$ on the graded vector space by this character.
Let us define the fine grading as $${}_{(q_1,.,q_d)}F_1(l_1,.,l_{d-p},\infty,.,\infty)_{l}:={}_{(q_1,.,q_d)}G_{1}(l_1,.,l_{d-p},0.,0)_{l}.$$\par
Consider $Q=(q_1,.,q_d)\in B_{\sigma_1}(\mathbb{T}),$ %and the set $$S_{\sigma_1}=\bigl\{Q':=(q_1,.,q_{d-p},q_{d-p+1}^{'},.,q_{d}^{'})\in B_{\sigma_1}(\mathbb{T})|(Q-Q')\in \frac{(\chi_{d-p+1},.,\chi_{d})\mathbb{Z}^d}{(m_{d-p+1},.,m_{d})}\bigl\}.$$ Note that $m_{d-p+1},.,m_d$ is a linear combination of $\chi_i,$ for $i=d-p+1(1)d.$ \newline
and let us describe $\hat{F}_{1,12}$ at the point $\sum_{i=1}^{d}(q_i+l_i)m^1_i.$\newline
%$$\bigoplus_{Q'\in S_{\sigma_1}}{}_{Q'}F_1(l_1,.l_{d-p},\infty,.,\infty).$$
Using the description of the graded tensor products, an element of $\hat{F}_{1,12}$ at the said point can be uniquely described as $$\bigoplus_{Q'\in S_{\sigma_1,Q}}(v_{Q'}\otimes \prod_{i=1}^{p}\lambda_i^{a_i})\ $$ where $$v_{Q'}\in {}_{Q'}G_{1}(l_1,.,l_d)$$ and % 
the homogeneous elements of the summands are elements of the module $$G_1\otimes_{\mathbb{C}[x_{\sigma_{1}(1)},.,x_{\sigma_{1}(d)}]} \mathbb{C}\bigg[\tau_1,.,\tau_{d-p},\lambda_{1}^{\pm 1},.,\lambda_{p}^{\pm 1}\bigg] .$$\par
The above summation can be explained by noting that the indexing set of the summation is given by
$$S_{\sigma_1,Q}:=\big \{Q'\in B_{\sigma_{1}}(\mathbb{T})\bigg|Q'=(Q-\sum_{i=1}^{p}a_{i}\chi_{d-p+i})  ,\ \prod_{i=1}^{p}\lambda_{i}^{a_i}\in B_1  \big \}.$$
Given an element of $B_1$, given by $(a_1,..,a_p)$ the corresponding $Q'$ is obtained by considering the box decomposition of the integer $Q-\sum_{i=1}^{p}a_{i}\chi_{d-p+i}$ with respect to the non-singular matrix given by $(m^1_1,..m^1_d).$\par For any two such choice of basis elements $\{a_i\},\{a'_i\},i=1,..,p$ corresponding to the same $Q'$ one can show that $$\begin{pmatrix}
  a_{1}-a'_{1} \\
  \vdots \\
  a_{p}-a'_p \\
\end{pmatrix}=[C_{p\times p}]\begin{pmatrix}
 l_1\\
 \vdots\\
 l_p\\
\end{pmatrix}$$ 
where $l_i \in \mathbb{Z}.$ Hence $\{a_i\},\{a'_i\}$ correspond to the same elements in the basis and hence the indexing set of the summand is in bijection with $|S_{\sigma_1,Q}|=|B_1|=|DG_{\sigma_2}/DH|.$

Let us compute the fine grading of the pullback module as a $K$ representation. Assuming $v_{Q'}$ is homogeneous of weight $l\in DG_{\sigma_1},$ one obtains the $X(K)$ weight of the element $v_{Q'}$ given by $\overline{l},$ the image of $l\in X(G_{\sigma_1})$ under the natural surjective map from $X(G_{\sigma_1})$ to $X(K)$ via the embedding into $X(G_{\sigma_1}\times G_{\sigma_2})$ (in [\ref{cokernel}]). Thus we obtain the $X(K)$ grading of $(v_{Q'}\otimes \prod_{i=1}^{p}\lambda_{i}^{a_i})$ given by $$(\overline{l}+\sum_{i=1}^{p}a_{i}\mu_{d-p+i}).$$ Hence, one obtains the value of $\hat{F}_{1,12}$ at the point $$\sum_{i=1}^{d}(q_i+l_i)m^1_i\ \in X(\mathbb{T})$$ given as a $X(K)$ representation by:
$$\bigoplus_{l\in DG_{\sigma_1}}\bigoplus_{\quad Q'\in S_{\sigma_1,Q}}{}_{Q'}G_{1}(l_1,.,l_d)_{l}\otimes {\sum_{i=1}^{p}a_i\mu_{d-p+i}}.$$ The above can be rewritten as:
$$\bigoplus_{\begin{array}{cc}
     &  l\in DG_{\sigma_1}\\
     & Q'\in S_{\sigma_1,Q}
\end{array}} \big ({}_{Q'}F_{1}(l_1,.,l_{d-p},\infty,.,\infty)_{l} \otimes_{i=1}^{p}\eta_{d-p+i}^{l_{d-p+i}}\ \bigl)\otimes \sum_{i=1}^{p}a_i\mu_{d-p+i}.$$\par
Consider $\mathscr{F}$ and restrict it to $Z_2$ and let us compute the value of $\hat{F}_{2,12}$ at the point $\sum_{i=1}^{d}(p_i+l'_i)m^2_{i}.$\newline
Consider $P=(p_1,.,p_d)\in B_{\sigma_2}(\mathbb{T})$ and consider the set $$\bigl\{Q'\in B_{\sigma_2}(\mathbb{T})|Q'=(P-\sum_{i=1}^{p}a'_i\chi_{d-p+i}),\prod_{i=1}^{p}\lambda_{i}^{a_{i}'}\in B_{2}\bigl\}.$$ Thus, applying similar techniques as before, the value at the said point is given by:
$$\bigoplus_{{\begin{array}{cc}
     &  l\in DG_{\sigma_2}\\
     & Q'\in S_{\sigma_2,Q}
\end{array}}}\big ({}_{Q'}F_{2}(l_1^{'},.,l_{d-p}^{'},\infty,.,\infty)_{l}\otimes _{i=1}^{p}(\eta'_{d-p+i})^{l_{d-p+i}'}\big )\otimes \sum_{i=1}^{p}a_i^{'}\mu_{d-p+i}.$$
\par 
We denote the equivalence of categories up to isomorphism between the finely graded $S-$ families and the isomorphism is taken as identity and the Gluing Formula is given at any point as the equality of $K$ representations determined as above. \par
From the description of $\hat{F}_{2,12}$ we see that it is enough to compute the Gluing formula at the point $$Q:=[\chi_1,..,\chi_{d-p}\  \chi_{d-p+1},..,\chi_{d}]\begin{pmatrix}
  q_1+l_1 \\
  \vdots \\
  q_{d-p}+l_{d-p} \\
  q_{d-p+1} \\
  \vdots \\
  q_{d} \\
\end{pmatrix}=\sum_{i=1}^{d-p}(q_i+l_i)\chi_i+\sum_{i=d-p+1}^{d}q_i\chi_i.$$ Here $(q_1,..,q_d)\in \text{Box}(\chi_1,..,\chi_d),\ l_i\in \mathbb{Z}.$\newline
Thus for $\hat{F}_{1,12}$ enough to evaluate at the point given by,
$$[C]^{-1}\begin{pmatrix}
  q_1+l_1 \\
  \vdots \\
  q_{d-p}+l_{d-p} \\
  q_{d-p+1} \\
  \vdots \\
  q_{d} \\
\end{pmatrix}=\left(\begin{array}{cc} Id_{d-p \times d-p} & 0\\ -C^{-1}_{p\times p}C'_{p\times d-p} & C^{-1}_{p\times p} \end{array}\right)\begin{pmatrix}
  q_1+l_1 \\
  \vdots \\
  q_{d-p}+l_{d-p} \\
  q_{d-p+1} \\
  \vdots \\
  q_{d} \\
\end{pmatrix}$$ with respect to the basis $(m^1_1,..,m^1_d).$ \newline 
Under the box decomposition, the co-ordinates, say are given by $$(q_1,..,q_{d-p},t_{1},..,t_p)\in B_{\sigma_1}(\mathbb{T})$$ and $$(l_1,..,l_{d-p},l_{d-p+1},..,l_d)\in \mathbb{Z}^d.$$
Similarly, the same computation for the other top cone lets us compute $\hat{F}_{2,12}$, given by the co-ordinates
$$(q_1,..,q_{d-p},t'_1,..,t'_p)\in B_{\sigma_2}(\mathbb{T}),$$ and $$(l_1,..,l_{d-p},l'_{d-p+1},..,l'_{d})\in \mathbb{Z}^d$$ with respect to the basis $(m^2_1,..,m^2_d).$\par 
%For the rest of the proof we are going to make a technical assumption. We prove the Gluing formula for toric DM stacks where the associated simplicial fan satisfies a particular property that is for each intersection of top cones $\sigma_i$ and $\sigma_j$, for every $v\in B_{\sigma_i}(\mathbb{T})$ can be uniquely written as $v=v'+\chi(b_j),\ v'\in B_{\sigma_j}(\mathbb{T}),b_j\in B_j.$\par
%Note that a certain point $z\in X(\mathbb{T})$ can be uniquely expressed as
%$$z:=\sum_{i=1}^{d}(q^j_i+l^j_i)m^j_i%=\sum_{i=1}^{d}q_im_i+\sum_{i=1}^{d-p}m_i^{'}l_{i}+\sum_{i=1}^{d}\chi(\alpha_{\sigma_1(i)}(\lambda))l_{i}
%$$ corresponding to the $\mathbb{T}$ weights corresponding to the top cones where
%Also note that $$\sum_{i=1}^{d}q_i\chi(\alpha_{\sigma_{1}(i)}(\lambda))=\sum_{i=1}^{d}\sum_{j=1}^{p}q_it_{i,j}\mu_{j}=\sum_{j=1}^{p}\mu_{j}(\sum_{i=1}^{d}q_it_{i,j})$$ where,
%$$\alpha_{\sigma_{1}(i)}(\lambda)=\prod_{j=1}^{p}\lambda_{j}^{t_{i,j}}.$$
%Also, observe that the above expression in terms of the the linearly independent set $$\bigg(m'_1,.,m'_{d-p},\chi(\alpha_{\sigma_{2}(d-p+1)}(\lambda)),.,\chi(\alpha_{\sigma_{2}(d)}(\lambda))\bigg)$$ can be rewritten as,
%$$=\sum_{i=1}^{d}p_im'_i+\ \sum_{i=1}^{d-p}m'_il'_i+\sum_{i=d-p+1}^{d}m'_il'_i\ .$$ 
%$$\bigg \{m^j_i:=m(\rho^j_{i})\bigg \}_{j=1,.,k\ i=1,.,d}$$ denote the $\mathbb{T}$ weights of the rays corresponding to $\sigma_{j}(i)$
%$$\sum_{i=1}^{d}((q_i+l_i))m_i^{1}=\sum_{i=1}^{d}(p_i+ l'_i)m_{i}^{2}$$\newline
Then we have the following 
%The above discussion says $\Hat{F}_{1,12}$ at the point $z \in X(\mathbb{T})$ can be computed by taking the tensor product of %$\Hat{F}_{1,12}$ at the point $\sum_{i=1}^{d-p}(q_i+z_i)m_i+\sum_{i=d-p+1}^{d}(q_i+l_i)m_i$ with the $X(\mathbb{T})$ weight of the element %$$(\prod_{i=1}^{d-p}\prod_{r=1}^{p}\tau_i\lambda_{r}^{t_{i,r}})^{l_i-z_i}$$ where
%$$\chi(\prod_{j=1}^{p}\tau_i\lambda_r^{t_{i,r}})=m_i,\ \chi(\tau_i)=m'_i\ \ \forall i=1,.,d-p.$$
%Thus we obtain,
%$$\Hat{F}_{1,12}(z)=\Hat{F}_{2,12}(z)\otimes \sum_{i=1}^{d-p}(l'_i-l_i)\mu_{i}-\sum_{r=1}^{p}t_{i,r}\mu_{d-p+r},$$ with the final %expression being as:

\begin{equation}\label{eqn_Gluing}
\begin{aligned}
&\bigoplus_{\begin{array}{cc}
&  l\in DG_{\sigma_1}\\
& Q'\in S_{\sigma_1,Q}
\end{array}} ({}_{Q'}F_{1}(l_1,.,l_{d-p},\infty,.,\infty)_{l} \otimes\ \sum_{i=d-p+1}^{d}\ (l_i\eta_i,0)) \otimes \sum_{i=d-p+1}^{d}a_{i-d+p}\mu_{i} \\
=&\bigoplus_{{\begin{array}{cc}
     &  l\in DG_{\sigma_2}\\
     & Q'\in S_{\sigma_2,Q}
\end{array}}}({}_{Q'}F_{2}(l_1,.,l_{d-p},\infty,.,\infty)_{l}\otimes\  \sum_{i=d-p+1}^{d} \ (0,l'_i\eta'_i))\otimes\ \sum_{i=d-p+1}^{d}a'_{i-d+p}\mu_{i}.\\ %& ~~~~~~~~~~~~~~~~~~~~~~~~~~~~~~~~~~~~~~~~~~~~+ 
%\sum_{i=1}^{d}\overline{(0,(l'_{i}-l_{i})\eta'_{i}}).
\end{aligned}
\end{equation}
\end{proof}

\end{document}
